\begin{table*}[t]
\centering
\scriptsize
\setlength{\tabcolsep}{4pt}
\renewcommand{\arraystretch}{0.98}
\begin{tabularx}{\textwidth}{@{}p{0.18\textwidth}p{0.16\textwidth}p{0.21\textwidth}p{0.15\textwidth}X@{}}
\toprule
Route & Frozen evidence & Readout type & Question addressed & Key takeaway \\
\midrule
Feature-space exploration & multiscale SAM \texttt{backbone\_fpn} features & clustering-only probe over pooled coarse features, flip-averaging, scale controls & RQ2 & Frozen features already exhibit nontrivial texture organization; coarse scale and positionality matter. \\
Proposal-/mask-space exploration (RWTD) & generated proposal masks plus dense rescue candidates & compatibility reasoning, conservative component selection, selective rescue & RQ3 and part of RQ4 & Natural scenes often externalize texture as fragmented masks, so mask-level commitment matters. \\
Proposal-/mask-space exploration (STLD) & generated proposal masks & compatibility reasoning and component selection & RQ3 and part of RQ4 & Canonical synthetic scenes often already contain a strong singleton mask. \\
Supporting appendix routes & generated proposal masks on ControlNet bridge and CAID & same proposal-space readout family under route-specific protocols & scope / breadth checks & Useful for context, but not mixed into the main matched comparison blocks. \\
\bottomrule
\end{tabularx}
\caption{Two-route map for the paper. The feature and proposal-/mask-space routes inspect different frozen signals and therefore answer different parts of the central question. The feature route reveals latent texture organization in the representation; the proposal-/mask-space route reveals how much of that organization is already present in the generated masks and recoverable under matched evaluation.}
\label{tab:route_protocol}
\end{table*}
